% ============================================================================
% Fichier cible : ARIEL_LOGMINER_MEMOIRE_FINAL/chapters/chapitre6_discussion.tex
% Passage à remplacer : fichier entier, de « \chapter{Discussion générale} »
%   jusqu'à la fin.
% Agent : Agent 3 — État de l'art et positionnement, chantier 2026-11-05.
% Sources : littérature déjà citée au chapitre 1 et références proposées par
%   l'Agent 3. Les résultats numériques proviennent exclusivement du corps du mémoire.
% ============================================================================

\chapter{Discussion générale}

\section{Du score isolé à la validité du protocole}

Le résultat transversal le plus solide ne porte pas sur la supériorité d'un détecteur,
mais sur la dépendance d'un score à son unité d'analyse, sa partition et ses baselines.
Cette dépendance est connue en détection d'intrusion : Sommer et Paxson soulignent le
décalage entre le monde fermé du laboratoire et les réseaux ouverts
\cite{sommer2010outside}, tandis qu'Arp et al. recensent les biais de sélection, les
fuites et les comparaisons inadéquates qui fragilisent l'apprentissage appliqué à la
sécurité \cite{arp2022dos}. Les expériences de ce mémoire en fournissent une réplication
sur quatre jeux, et non une découverte théorique.

Sur CICIDS2017, une partition aléatoire produit un résultat très favorable, alors que le
holdout par scénario et la séparation temporelle exposent une faible généralisation. Ce
contraste rejoint les erreurs de génération, de caractéristiques et d'étiquetage déjà
documentées pour ce corpus \cite{engelen2021troubleshooting,lanvin2023errors}. Sur
CSE-CIC-IDS2018, le transfert entre deux journées dépend fortement du port de destination
et s'effondre lorsque le sens est inversé. Liu et al. ont indépendamment identifié, dans
CICIDS2017 et CSE-CIC-IDS2018, des erreurs couvrant l'orchestration, l'extraction et les
labels \cite{liu2022error}. Le résultat du mémoire doit donc être lu comme un audit local
de deux fichiers, pas comme une validation externe d'un détecteur d'intrusion.

HDFS illustre une autre menace : l'événement n'est pas l'unité annotée, le bloc l'est.
L'agrégation modifie alors la matrice de confusion et l'interprétation du F1. Sur le test
complet, la baseline de surprise moyenne des templates atteint un F1 de $0{,}872434$,
légèrement supérieur aux $0{,}868373$ d'Histogram : son avantage prédictif n'est donc pas
établi. Sur BGL, le score principal coïncide avec la règle « sévérité absente de l'entraînement normal » ; le
détecteur ne met en évidence aucune structure plus riche sous ce protocole. Les travaux
de Le et Zhang montrent précisément que la sélection des données, le regroupement et le
déséquilibre peuvent transformer l'évaluation des détecteurs de journaux
\cite{le2022logbased}. L'ensemble de ces observations impose de présenter chaque score
avec son unité, sa partition et la meilleure règle simple évaluée sur le même test.

La comparaison préenregistrée E-H3b entre modèles communs et spécialisés, sous
affectation parfaite de la famille, n'établit pas la non-infériorité pour
\texttt{linux\_auth} : la borne inférieure de l'IC apparié à 95~\% est
$-0{,}012326$, sous la marge de $-0{,}01$. Cette expérience ne mesure pas la qualité du
routeur et ne remplace pas une comparaison avec une table statique.

\section{Une intégration qui n'est pas une lacune fonctionnelle}

L'assemblage de l'ingestion, de la détection, de l'alerte et de la visualisation existe
déjà dans des cadres établis. Wazuh documente une architecture composée d'agents, d'un
serveur d'analyse, d'un indexeur et d'un tableau de bord
\cite{wazuh2026architecture}. Elastic Security associe données normalisées, règles,
corrélation et investigation \cite{elastic2026ecs,elastic2026security}. OpenSearch
Security Analytics fournit des détecteurs, des règles, des constats, des alertes et un
moteur de corrélation \cite{opensearch2026securityanalytics}. Dans la recherche, LogAI et
Loglizer offrent déjà des interfaces ou outils communs pour plusieurs méthodes d'analyse
de journaux \cite{cheng2023logai,he2016experience,logpai2026loglizer}.

Ariel Logminer ne comble donc pas un vide entre algorithme et plateforme. Il ne compare
d'ailleurs ni sa couverture, ni son débit, ni sa qualité de détection à ces cadres. Son
apport défendable est plus restreint : une chaîne locale où l'on peut relier des résultats
à des protocoles et isoler certains mécanismes. Cette instrumentation a une valeur
d'ingénierie et de pédagogie expérimentale. Elle devient une contribution scientifique
seulement lorsque l'expérience oppose effectivement des explications concurrentes. Les
audits de partitions et de raccourcis satisfont partiellement cette exigence ; la simple
présence de composants dans une architecture ne la satisfait pas.

\section{Contract Net face à l'affectation statique}

Contract Net est conçu pour allouer une tâche à partir d'offres produites par des
contractants \cite{smith1980contract,fipa2002contractnet}. Son intérêt théorique apparaît
lorsque les agents détiennent un état local pertinent, que les offres peuvent diverger et
que la négociation améliore une décision par rapport à une règle plus simple. Ces
conditions ne sont pas démontrées par la seule émission de messages
\texttt{CFP}, \texttt{PROPOSE} et \texttt{AWARD}.

Le concurrent pertinent est une file de tâches avec table de routage ou files nommées.
Celery affecte des tâches nommées à des workers ; Kafka et Redis répartissent des messages
entre consommateurs ; Temporal associe des files de tâches à des workers et reconstruit
l'état d'un workflow depuis son historique
\cite{celery2026tasks,kafka2026semantics,redis2026streams,temporal2026tasks}. Ces cadres
peuvent fournir découplage, tentatives, acquittements et reprise sans négociation
Contract Net.

Dans les campagnes présentées, les poids et les capacités sont configurés, une composante
de l'utilité est constante, deux composantes sont redondantes dans le contexte évalué et
le seuil de refus par faible utilité est inactif. La répartition de l'endurance suit les
profils fixés ; lors de la détection HDFS distribuée, les utilités égales sont départagées
par un ordre déterministe et toutes les tâches vont au même agent. Ces observations
rapprochent l'allocation d'une table statique. Elles ne réfutent pas l'intérêt général de
Contract Net, mais elles ne démontrent pas son bénéfice dans cette instanciation.

Le benchmark montre en outre un coût substantiel de l'orchestration sur les micro-tâches,
et la seconde machine n'accroît pas la capacité observée dans la campagne de saturation.
Une comparaison concluante doit conserver transport, workers, tâche et mécanismes de
reprise, puis changer uniquement la politique d'affectation. Sans cette ablation, la
conclusion doit rester descriptive : l'échange négocié fonctionne, mais son avantage sur
une affectation déterministe n'est pas établi.

\section{Reprise, livraison et idempotence}

La reprise observée après l'arrêt d'un worker combine deux mécanismes. Le transport
conserve ou réattribue une tâche non acquittée ; le magasin d'idempotence empêche que la
nouvelle tentative répète l'effet métier. Redis Streams documente une sémantique au moins
une fois pour les groupes de consommateurs : les messages non acquittés restent en
attente et peuvent être réclamés par un autre consommateur
\cite{redis2026streams}. Celery avertit de même qu'un acquittement tardif peut conduire à
plusieurs exécutions et exige des tâches idempotentes \cite{celery2026tasks}. Kafka
rappelle que l'exactement-une-fois dépend du périmètre transactionnel considéré
\cite{kafka2026semantics}.

L'ablation appariée du mémoire confirme cette séparation : l'ajout du même magasin aux
architectures de référence supprime aussi les effets dupliqués. L'idempotence n'est donc
pas un bénéfice propre aux agents ou à Contract Net. La contribution consiste à l'avoir
intégrée et testée dans la chaîne, pas à en avoir montré la spécificité architecturale.
La reprise multi-VM établit un scénario précis d'arrêt après persistance et avant
publication du résultat. Elle ne couvre pas une panne de Redis, une partition réseau ou
un basculement de stockage ; ces propriétés requièrent des protocoles distincts.

\section{Hétérogénéité, normalisation et corrélation}

La lecture de plusieurs formats démontre une couverture d'adaptateurs, non une sémantique
commune complète. ECS spécifie des champs, des types et des conventions d'extension pour
normaliser des événements \cite{elastic2026ecs}. OCSF fournit une taxonomie et des classes
d'événements indépendantes du fournisseur et du stockage \cite{ocsf2026schema}. Le
contrat d'Ariel Logminer reste un schéma local plus réduit. Le fait que les flux
tabulaires ne renseignent aucun des sept champs standard mesurés interdit de conclure à
une normalisation homogène. De plus, chaque source publique est évaluée séparément : le
mémoire ne présente aucun incident réel analysé conjointement à partir d'un journal
système et d'un flux réseau.

La corrélation met en évidence la même différence entre intégration fonctionnelle et
validation scientifique. Les cadres classiques distinguent fusion d'alertes similaires,
corrélation de scénarios, vérification et priorisation
\cite{cuppens2002alert,valeur2004correlation}. Le mécanisme étudié ici est une règle de
regroupement par fenêtre et égalité de clés, sans modèle causal. Sur le corpus synthétique
construit avec ces clés, il obtient une forte précision, mais son rappel diminue et une
baseline « hôte + fenêtre » le dépasse lorsque les incidents enchaînent plusieurs
catégories. Cette défaite est informative : ajouter des clés réduit les fusions au prix
de fragmenter précisément les scénarios multi-étapes. Une validation externe doit porter
sur des incidents authentifiés, non utilisés pour concevoir les règles.

\section{Mémoire, gouvernance et rôle de l'analyste}

La mémoire persistante assure une traçabilité et fournit un état aux décisions. Elle ne
démontre pas, par sa présence, une adaptation bénéfique. Les différences appariées ne
mettent pas en évidence de gain systématique de débit, de latence ou d'équité ; le statut
de la mémoire doit donc être celui d'une infrastructure réalisée, dont l'effet reste à
évaluer selon une métrique préalablement choisie.

La littérature MLOps étend cette exigence à tout le cycle de vie. Les dépendances de
données, les boucles de rétroaction et la configuration créent une dette technique propre
aux systèmes d'apprentissage \cite{sculley2015debt}. La préparation à la production
requiert des tests et une surveillance portant sur les données, le modèle et
l'infrastructure \cite{breck2017mltest}. Les études de terrain soulignent aussi la
difficulté de versionner les données et de maintenir des modèles liés à leur contexte
logiciel \cite{amershi2019software}. Enfin, la gouvernance suppose des rôles, une
documentation, une mesure et une revue continues \cite{nist2023airmf}.

Le registre, le mode \textit{dry-run} et les chemins de promotion du prototype sont donc
des briques, non une validation MLOps. Aucun cycle répété ne relie candidat, décision,
déploiement, surveillance et retour arrière. De même, le tableau de bord offre une
restitution fonctionnelle, mais aucune étude avec des analystes ne mesure le temps de
triage, l'exactitude ou la charge cognitive. L'exploitabilité humaine et l'amélioration
produite par la gouvernance restent deux questions ouvertes.

\section{Nature et limites de la contribution}

La contribution de ce travail est une intégration instrumentée accompagnée d'une étude
critique de ses protocoles. Trois apports sont établis. Premièrement, le prototype relie
plusieurs familles de sources à des traitements explicites et conserve des artefacts de
traçabilité. Deuxièmement, les campagnes exposent la fragilité de scores favorables face
à des partitions, unités et règles triviales mieux choisies. Troisièmement, les ablations
séparent certains mécanismes : l'idempotence relève du magasin, la reprise de la
combinaison transport--application, et le coût de l'orchestration n'est pas compensé dans
les charges étudiées.

Ces apports ne constituent ni un nouvel algorithme, ni un nouveau schéma, ni une preuve de
supériorité multi-agent. Ils répliquent plusieurs critiques publiées et les illustrent
dans une chaîne locale. La validité externe est limitée par des jeux publics construits,
deux fichiers CSE-CIC-IDS2018 dont l'intégrité complète n'est pas établie, une fixture
Apache minimale, des corpus synthétiques de corrélation et une distribution de décision
où l'inférence n'est pas recalculée sur les machines virtuelles. La validité de construit
est limitée lorsque « normalisé », « autonome » ou « distribué » désigne l'existence d'un
mécanisme plutôt que le bénéfice qu'il produit.

Le positionnement final est donc celui d'un prototype d'ingénierie expérimentale qui
documente autant ses échecs que ses réussites. Sa valeur scientifique réside dans les
comparaisons qui excluent une explication simple ; là où le comparateur manque, la
revendication reste non tranchée.

\section{Priorités de prolongement}

Les prolongements doivent cibler les menaces qui empêchent actuellement de trancher les
revendications. La première priorité est de comparer Contract Net à une affectation
statique sur les mêmes tâches et le même transport. La deuxième est d'évaluer un scénario
réellement hétérogène où un même incident possède des traces système et réseau, mappées à
un schéma explicite tel qu'ECS ou OCSF. La troisième est de valider la corrélation sur des
incidents multi-étapes externes au générateur.

Sur le plan prédictif, les prochaines campagnes doivent utiliser des données corrigées ou
réextraites, des partitions externes, des baselines au même niveau et des règles de seuil
fixées avant l'exécution. Sur le plan opérationnel, il faut distribuer le calcul complet,
tester des pannes du broker et des partitions réseau, puis mesurer les charges où le coût
fixe de l'orchestration pourrait devenir secondaire. Enfin, un cycle de gouvernance doit
versionner données, parseur, caractéristiques et modèle, puis tester promotion,
surveillance et retour arrière. Ces expériences ne garantissent pas un résultat favorable ;
elles déterminent si l'intégration produit une connaissance qui dépasse son propre
fonctionnement.
